%% Stage report II (IEEEtran, English) — independent prospective cohort (C1): lock, execution, verification
%% Compile: pdflatex -interaction=nonstopmode stage_report_02_ieee_en.tex  (twice)
%% Companion: stage_report_02_ieee_zh.tex (independent, xelatex + xeCJK)
\documentclass[conference]{IEEEtran}

\usepackage{amsmath}
\usepackage{amssymb}
\usepackage{booktabs}
\usepackage{graphicx}
\usepackage{float}
\usepackage[T1]{fontenc}
\usepackage{lmodern}
\usepackage{microtype}
\usepackage{url}

\begin{document}

\title{Stage Report II: Independent Prospective Cohort (C1)\\\ Lock, Execution and Verification Status}

\author{\IEEEauthorblockN{Dongxu Jiang}
\IEEEauthorblockA{Email: jiangdongxu04@gmail.com}}

\maketitle

\begin{abstract}
This report records the \emph{execution status} of the independent prospective cohort study (C1)
as of 2026-09-14: a locked twenty-asset, single-provider cohort; a predictor frozen before any
outcome exists; a 120-run panel delegated to a peer machine that is offline-capable and idempotent;
and the intake protocol that accepts or rejects each returned batch. It deliberately contains
\textbf{no analysis of C1 outcomes}.
\end{abstract}

\begin{IEEEkeywords}
prospective validation, pre-registration, blinded analysis, cohort locking, execution provenance,
intake verification.
\end{IEEEkeywords}

\section{Purpose and Scope}

The programme has moved from exploratory diagnostics to a \emph{prospective} test. This report
documents the state of that test's \emph{machinery} and nothing else.

\section{Programme Sequence and Current Position}
\label{sec:seq}

Four stages have been executed in order, each gated by the previous one:

\begin{table}[!t]
\centering
\footnotesize
\begin{tabular}{@{}cll@{}}
\toprule
Stage & Content & Status \\
\midrule
S1 & Reproduction and controlled diagnostics & complete \\
S2 & Structural validation of the candidate regularity & \textbf{FAIL} \\
S3 & Capacity-controlled breadth panel (F1) & \textbf{SUCCESS} \\
S4 & Independent prospective cohort (C1) & executing \\
\bottomrule
\end{tabular}
\caption{Programme sequence. Each stage is gated by the previous one; S4 is the subject of this
report.}
\label{tab:seq}
\end{table}

S1 established the reproduction and the operator-specialisation diagnostics. S2 tested the candidate
structural regularity on four pre-registered assets and returned a failure, triggered by a capacity
contradiction: an asset-level operator preference whose sign is not stable across model widths. S3
then asked the more basic question and answered it in the narrow sense that asset-level preference
replicates on a majority of a seven-asset panel while remaining a \emph{tendency}, not a property of
an asset. S4 --- the present stage --- tests the candidate association \emph{prospectively} on a
genuinely independent cohort.

\section{The Locked Cohort}
\label{sec:cohort}

The cohort was locked before any outcome was inspected, and the lock is a table of hashes rather than
a list of intentions.

\begin{table}[!t]
\centering
\footnotesize
\begin{tabular}{@{}lrl@{}}
\toprule
Property & Value & Note \\
\midrule
Assets & 20 & pre-named universe, no additions \\
Provider & Yahoo & one provider only; no composite \\
Window & 2016-01-04/01 -- 2025-12-31 & daily OHLC, unadjusted \\
Exclusions & 0 & all 20 delivered and accepted \\
Rows per asset & 2444--2602 & market calendar differences \\
Hash table & per file & verified on intake \\
\bottomrule
\end{tabular}
\caption{Cohort lock summary. The per-file SHA-256 table is the authoritative record; row counts
differ across assets because exchange calendars differ, not because data were imputed.}
\label{tab:cohort}
\end{table}

Two properties matter for what follows. The cohort shares the \emph{provider} of the discovery set
while differing in \emph{market}: its assets exclude both the original ten and the four assets used in
S2. Rows were dropped only where the provider returned all-null price rows and were never
interpolated, with the drop counts documented per file. Acquisition was performed by the peer machine,
because the analysing host cannot reach the provider.

\section{The Frozen Predictor}
\label{sec:pred}

The predictor is frozen \emph{before} any outcome exists, which is the difference between a
prospective test and a descriptive one. Its construction rule:

\begin{itemize}
\item a feature window covering closes $[i, i+96)$ is usable only if the whole window ends at or
before the decision point, i.e.\ $i+96 \le \mathrm{border1\_test}$;
\item the decision point follows the frozen split convention,
$\mathrm{border1\_test} = n - \lfloor 0.2n \rfloor - 96$ with $\mathrm{num\_train} = \lfloor 0.7n
\rfloor$;
\item the asset-level predictor aggregates the window features over the usable windows only, so the
last close index touched is $\mathrm{border1\_test} - 1$.
\end{itemize}

A leakage check is executed and printed: it asserts that every index used by the predictor lies
strictly below the decision point. The check passed for all twenty assets. A contemporaneous
analogue computed over the test windows is carried as a \emph{contrast only} and is never used for
the primary test. Because the predictor is frozen, it may not be recomputed with a different window,
estimator or cutoff once outcomes are seen.

\section{Delegated Execution}
\label{sec:exec}

S4 requires 120 training runs ($20$ assets $\times$ $2$ arms $\times$ $3$ seeds) under the unchanged
protocol, with per-run predictions and ground truth retained for the paired analysis.
\begin{table}[!t]
\centering
\footnotesize
\begin{tabular}{@{}lrl@{}}
\toprule
Property & Value & Note \\
\midrule
Runs & 120 & $20 \times 2 \times 3$ \\
Concurrency & 2 lanes & frozen; widening declined \\
Threads per process & 14 & unchanged mid-campaign \\
Retained per run & metrics, pred, true & command line and log tail too \\
Checkpoints & not transferred & not needed for the analysis \\
Batch reporting & 40 runs or 6 h & manifest per batch \\
Progress at writing & 0/120 & launched 13:36, first runs $\approx$60 min \\
\bottomrule
\end{tabular}
\caption{Execution configuration. Throughput changes that would alter run comparability mid-campaign
(thread count, loader workers, extra lanes) were declined and deferred.}
\label{tab:exec}
\end{table}

The executor was selected because its network reaches the provider while the analysing host does not.
Its runner is \emph{offline-capable} by construction: everything is staged once, after which training
reads only local files, so a network outage cannot interrupt the campaign. Resilience rests on four
mechanisms: idempotent skipping of completed runs; a watchdog that restarts lanes without ever
duplicating a run; detection of stalled runs by \emph{process-tree CPU time} rather than log output;
and a contingency plan that pre-authorises self-recovery while forbidding any protocol or data
change. One genuine defect was found and fixed during the campaign: a restart had orphaned two
finished runs that no lane recorded, which would have produced a silent two-run gap; the runner now
\emph{ingests} such runs and marks them in the manifest, so the process, not merely the content, is
auditable.

\section{Intake Verification}
\label{sec:verify}

Nothing enters the evidence chain unverified. Each returned batch is checked mechanically before any
file is accepted:
\begin{enumerate}
\item \textbf{Protocol conformance per run:} the stored command line must contain every frozen
hyperparameter value together with that run's own seed, model and data path.
\item \textbf{Content integrity:} predictions and ground truth must match the batch manifest hashes.
\item \textbf{Metric integrity:} the metric array must parse as the frozen six-value order.
\item \textbf{Whole-batch rejection:} if any run, manifest row, hash or command line disagrees, the
\emph{entire} batch is rejected and returned with a difference list; partial results are never
copied. The verifier's default mode is a dry run.
\item \textbf{Handoff inventory:} an accepted batch yields a per-file SHA-256 inventory and a handoff
receipt recording the frozen revision actually executed.
\end{enumerate}

One gap cannot be closed by the receiving side alone: transferred files can be shown to match the
batch manifest, but only the executor can show that its local originals match what was transferred.
That confirmation is therefore a required handoff item, not an optional courtesy.

\section{Independence and Governance}
\label{sec:gov}

Three roles are kept separate: the \emph{executor} runs the panel and returns raw artefacts without
performing any analysis; the \emph{analyst} computes the predictor and the headline statistic; a
\emph{blind third party}, which has seen neither the analysis nor its own expectations, recomputes
the headline statistic from the raw artefacts and reports a difference list, with the explicit rule
that a discrepancy is a finding and never something to be tuned away. Executor and analyst are both
excluded from that third role, and readers of this report are asked not to treat it as evidence
about C1's outcome.

\section{Milestones and Status}
\label{sec:mil}

\begin{table}[!t]
\centering
\footnotesize
\begin{tabular}{@{}rlp{0.52\columnwidth}@{}}
\toprule
\# & Milestone & Status \\
\midrule
M1 & Cohort acquired by the peer machine & done \\
M2 & Intake validation of the cohort & done, 20/20 \\
M3 & Cohort locked with a per-file hash table & done \\
M4 & Predictor frozen with a leakage check & done, all assets \\
M5 & Verifier hardened and smoke-tested & done \\
M6 & Panel launched (120 runs, 2 lanes) & running, 0/120 \\
M7 & First batch returned and verified & pending \\
M8 & Blind third-party recomputation & pending \\
\bottomrule
\end{tabular}
\caption{Milestones. M7 gates M8, and no stage beyond M7 is attempted before the panel completes.}
\label{tab:mil}
\end{table}

\section{Next Steps}
\label{sec:next}

\begin{enumerate}
\item Continue execution and health monitoring; keep the campaign's conditions constant.
\item Verify each returned batch under the intake protocol; reject whole batches on any mismatch.
\item On completion, assemble the handoff: runs, prediction and ground-truth arrays, the batch
manifest, the frozen revision actually executed, the per-file hash inventory, and the executor's
local-origin hash confirmation.
\item Release the artefact set to the blind third party and await its difference list; only then read
the pre-committed interpretation of whichever branch the outcome falls into.
\end{enumerate}

\section{What This Report Deliberately Omits}
\label{sec:omit}

No estimate of the headline statistic appears here, no per-asset interaction effect, no reading of
either branch of the pre-committed interpretation, and no claim about the eligibility of any
structured operator. Those omissions are not gaps: they are the condition under which the report is
useful, and the same applies to any progress summary written while the panel runs.

\section*{Reproducibility Notes}
\addcontentsline{toc}{section}{Reproducibility Notes}

All artefacts cited here live in the project repository: the lock table with per-file hashes, the
frozen predictor, the intake verifier with its inventory and handoff receipt, and the batch manifests
that reproduce each run's command line.

\end{document}
